% Chapter 1, Topic from _Linear Algebra_ Jim Hefferon
%  http://joshua.smcvt.edu/linalg.html
%  2001-Jun-09
\topic{Analisis Masukan-Keluaran}
\index{Analisis Masukan-Keluaran|(}
Perekonomian adalah jejaring kesalingbergantungan yang luar biasa rumit.
Perubahan di satu bagian dapat merambat dan memengaruhi bagian-bagian lain.
Para ekonom telah lama berupaya agar dapat menggambarkan dan membuat
prediksi tentang objek serumit itu.
Model matematika yang menggunakan sistem persamaan linear
merupakan salah satu perangkat utama.
Salah satu contohnya adalah Analisis Masukan-Keluaran, yang dirintis oleh
W.~Leontief,\index{Leontief, W.}
peraih Penghargaan Nobel Ekonomi tahun 1973.

Perhatikan suatu perekonomian yang terdiri atas banyak bagian, dua di antaranya
adalah industri baja dan industri otomotif.
Keduanya saling berinteraksi erat
ketika berusaha memenuhi permintaan atas produk mereka
dari bagian-bagian lain perekonomian, yakni pengguna di luar sektor baja dan
otomotif.
Sebagai contoh, jika permintaan eksternal terhadap mobil meningkat, hal itu
akan meningkatkan penggunaan baja oleh industri otomotif.
Atau, jika permintaan eksternal terhadap baja menurun, industri baja akan
mengurangi pembelian truk.
Model yang
kita bahas di sini menerima permintaan eksternal sebagai masukan
dan memprediksi cara kedua industri itu berinteraksi untuk memenuhinya.

Kita mulai dengan statistik produksi dan konsumsi.
(Angka-angka ini, yang menyatakan nilai dalam jutaan
dolar, berasal dari
\cite{Leontief1965} yang menggambarkan perekonomian A.S.\ tahun 1958.
Statistik masa kini akan berbeda karena
inflasi dan perubahan teknologi di kedua industri tersebut.)
\begin{center}
  \begin{tabular}{r|rrrr}
         \multicolumn{1}{c}{\ }  %put the | in the right place
         &\multicolumn{1}{r}{\begin{tabular}[b]{@{}c@{}}
                                \textit{digunakan oleh} \\[-.65ex] \textit{baja}
                             \end{tabular}}
         &\multicolumn{1}{r}{\begin{tabular}[b]{@{}c@{}}
                                \textit{digunakan oleh} \\[-.65ex] \textit{otomotif}
                             \end{tabular}}
         &\multicolumn{1}{r}{\begin{tabular}[b]{@{}c@{}}
                                \textit{digunakan oleh} \\[-.65ex] \textit{lainnya}
                              \end{tabular}}
         &\textit{total}                                                \\
         \cline{2-5}
    \begin{tabular}{r} \textit{nilai} \\[-.65ex] \textit{baja} \end{tabular}
         &$5\,395$  &$2\,664$  &     &$25\,448$                          \\
    \begin{tabular}{r} \textit{nilai} \\[-.65ex] \textit{mobil} \end{tabular}
         &$48$      &$9\,030$  &     &$30\,346$
  \end{tabular}
\end{center}
Sebagai contoh, nilai dolar baja yang digunakan oleh industri otomotif pada
tahun ini adalah $2,664$ juta.
Perhatikan bahwa suatu industri dapat mengonsumsi sebagian hasil produksinya sendiri.

Kita dapat melengkapi permintaan eksternalnya.
Nilai baja yang digunakan pihak lain tahun ini adalah $17,389$,
sedangkan permintaan eksternal terhadap mobil adalah $21,268$.
Dengan demikian, kita memperoleh gambaran lengkap mengenai
cara industri otomotif dan baja berinteraksi untuk memenuhi permintaan mereka.

Sekarang bayangkan bahwa permintaan eksternal terhadap baja belakangan ini
meningkat sebesar $200$ per tahun, sehingga kita memperkirakan nilainya tahun
depan akan menjadi $17,589$.
Kita juga memperkirakan bahwa permintaan eksternal terhadap mobil tahun depan
akan turun sebesar $25$ menjadi $21,243$.
Kita ingin memprediksi keluaran total tahun depan.

Prediksi tersebut tidak sesederhana menambahkan $200$ pada total baja tahun ini
dan mengurangkan $25$ dari total mobil tahun ini.
Salah satu alasannya, kenaikan produksi baja akan membuat
industri tersebut meningkatkan permintaannya terhadap mobil, sehingga
meredam penurunan permintaan eksternal terhadap mobil.
Di sisi lain, turunnya permintaan eksternal terhadap mobil akan membuat
industri otomotif menggunakan lebih sedikit baja sehingga sedikit mengurangi
peningkatan kegiatan usaha industri baja.
Singkatnya, kedua industri ini membentuk sebuah sistem.
Kita harus memprediksi titik keseimbangan sistem secara keseluruhan.

Berikut persamaannya.
\begin{align*}
  \text{produksi baja tahun depan}
     &=\text{penggunaan baja oleh industri baja tahun depan}   \\[-.8ex]
     &\quad\hbox{}+\text{penggunaan baja oleh industri otomotif tahun depan}  \\[-.8ex]
     &\quad\hbox{}+\text{penggunaan baja oleh pihak lain tahun depan} \\
  \text{produksi mobil tahun depan}
     &=\text{penggunaan mobil oleh industri baja tahun depan}   \\[-.8ex]
     &\quad\hbox{}+\text{penggunaan mobil oleh industri otomotif tahun depan}  \\[-.8ex]
     &\quad\hbox{}+\text{penggunaan mobil oleh pihak lain tahun depan}
\end{align*}
Untuk ruas kiri, misalkan
$s$ adalah produksi total baja tahun depan dan
$a$ adalah keluaran total mobil tahun depan.
Untuk ruas kanan, sebagaimana dibahas di atas,
perkiraan permintaan eksternal kita adalah $17,589$ dan $21,243$.

Untuk penggunaan baja oleh industri baja tahun depan, perhatikan bahwa tahun ini
industri baja menggunakan $5,395$ unit masukan baja untuk menghasilkan $25,448$
unit keluaran baja.
Jadi, tahun depan, ketika industri baja menghasilkan $s$ unit keluaran,
kita memperkirakan bahwa untuk melakukannya dibutuhkan
$s\cdot (5\,395)/(25\,448)$
unit masukan baja\Dash ini semata-mata merupakan asumsi bahwa masukan
sebanding dengan keluaran.
(Kita mengasumsikan bahwa rasio masukan terhadap keluaran tetap konstan dari waktu
ke waktu; dalam praktiknya, model dapat mencoba memperhitungkan kecenderungan
perubahan rasio tersebut.)

Penggunaan baja oleh industri otomotif tahun depan dihitung dengan cara serupa.
Tahun ini industri otomotif menggunakan $2,664$ unit masukan baja untuk
menghasilkan $30,346$ unit keluaran mobil.
Jadi, tahun depan, ketika keluaran total industri otomotif adalah $a$, kita
memperkirakan industri itu akan mengonsumsi
$a\cdot (2\,664)/(30\,346)$ unit baja.

Dengan melengkapi persamaan lainnya menggunakan cara yang sama, kita memperoleh
sistem persamaan linear berikut.
\begin{equation*}
  \begin{linsys}{3}
    (5\,395/25\,448)\cdot s
      &+ &(2\,664/30\,346)\cdot a &+ &17\,589
         &= &s \\
    (48/25\,448)\cdot s
      &+ &(9\,030/30\,346)\cdot a &+ &21\,243
         &= &a
  \end{linsys}
\end{equation*}
Dengan memindahkan variabel ke satu ruas dan konstanta ke ruas lainnya,
\begin{equation*}
  \begin{linsys}{2}
      (20\,053/25\,448)s &- &(2\,664/30\,346)a &= &17\,589 \\
     -(48/25\,448)s      &+ &(21\,316/30\,346)a &= &21\,243
  \end{linsys}
\end{equation*}
kemudian menerapkan Metode Gauss atau menggunakan \textit{Sage}, seperti berikut,
\begin{lstlisting}
sage: var('a,s')
(a, s)
sage: eqns=[(20053/25448)*s - (2664/30346)*a == 17589,  
....:        (-48/25448)*s + (21316/30346)*a == 21243] 
sage: solve(eqns, s, a)
[[s == (2745320544312/106830469), a == (6476293881123/213660938)]]
sage: n(2745320544312/106830469)
25697.9171767186
sage: n(6476293881123/213660938)
30311.0804517904
\end{lstlisting}
kita memperoleh prediksi:
$s=25,698$ dan $a=30,311$.

Di atas telah kita bahas bahwa memprediksi total tahun depan
tidak sesederhana menambahkan $200$ pada total baja tahun lalu dan
mengurangkan $25$ dari total mobil tahun lalu.
Perbandingan prediksi ini
dengan angka tahun berjalan
menunjukkan bahwa produksi total industri baja seharusnya naik sebesar~$250$,
sementara
total mobil turun sebesar $35$.
Kenaikan permintaan eksternal terhadap baja menimbulkan kenaikan
permintaan antarindustri oleh industri baja, yang sedikit diredam oleh
penurunan mobil, tetapi menghasilkan total yang lebih tinggi daripada
$200$.
Demikian pula, total mobil turun lebih dari~$25$ meskipun terbantu oleh
industri baja.

Salah satu keunggulan model matematika adalah kita dapat
mengajukan pertanyaan ``Bagaimana jika \ldots?''
Sebagai contoh, kita dapat bertanya,
``Bagaimana jika perkiraan permintaan eksternal tahun depan agak meleset?''
Untuk memahami seberapa besar perubahan prediksi model sebagai
reaksi terhadap perubahan perkiraan kita, kita dapat merevisi perkiraan
permintaan eksternal baja tahun depan dari $17,589$ turun menjadi
$17,489$, sambil mempertahankan asumsi bahwa permintaan eksternal terhadap
mobil tahun depan tetap sebesar $21,243$.
Sistem yang dihasilkan
\begin{equation*}
  \begin{linsys}{2}
        (20\,053/25\,448)s &- &(2\,664/30\,346)a &= &17\,489 \\
      -(48/25\,448)s      &+ &(21\,316/30\,346)a &= &21\,243
  \end{linsys}
\end{equation*}
memberikan $s=25,571$ dan $a=30,311$.
% sage: var('a,s')                                        
% (a, s)
% sage: eqns=[(20053/25448)*s - (2664/30346)*a == 17489, 
% ....:        (-48/25448)*s + (21316/30346)*a == 21243] 
% sage: solve(eqns, s, a)                             
% [[s == (2731759305112/106830469), a == (6476221050723/213660938)]]
% sage: n(2731759305112/106830469)                    
% 25570.9754968126
% sage: n(6476221050723/213660938)
% 30310.7395827449
Inilah yang disebut \definend{analisis sensitivitas}.
Kita sedang melihat seberapa sensitif prediksi model kita terhadap
ketepatan asumsi.

Tentu saja, kita dapat membahas model yang lebih besar dan merinci interaksi
di antara lebih banyak sektor perekonomian;
model semacam itu biasanya diselesaikan dengan komputer.
Selain itu, satu model tentu tidak cocok untuk setiap kasus, dan
memastikan bahwa asumsi yang mendasari suatu model
masuk akal bagi prediksi tertentu memerlukan pertimbangan para ahli.
Namun, dengan memperhatikan berbagai batasan tersebut,
model ini terbukti dalam praktik sebagai perangkat yang bermanfaat dan akurat
untuk analisis ekonomi.
Untuk bacaan lebih lanjut, lihat \cite{Leontief1951} dan \cite{Leontief1965}.


\begin{exercises}
  \item
    Dengan menggunakan sistem baja-otomotif di atas, perkirakan
    produksi total tahun depan untuk kasus-kasus berikut.
    \begin{exparts}
      \partsitem Permintaan eksternal baja tahun depan naik $200$ dari
         tahun ini, sedangkan permintaan eksternal mobil tidak berubah.
      \partsitem Permintaan eksternal baja tahun depan naik $100$ dan
         permintaan eksternal mobil naik $200$.
      \partsitem Permintaan eksternal baja tahun depan naik $200$ dan
         permintaan eksternal mobil naik $200$.
    \end{exparts}
    \begin{answer}
      Jawaban ini diperoleh dari \textit{Sage}.
      \begin{exparts}
        \partsitem \Sage{} memperoleh $s\approx 25\,952$ dan $a\approx 30\,312$.
\begin{lstlisting}
sage: var('s,a')
(s, a)
sage: system = [(20053/25448)*s -  (2664/30346)*a == 17789,
....:             -(48/25448)*s + (21316/30346)*a == 21243]
sage: solve(system, s,a)
[[s == (2772443022712/106830469), a == (6476439541923/213660938)]]
sage: n(2772443022712/106830469), n(6476439541923/213660938)
(25951.8005365305, 30311.7621898814)          
\end{lstlisting}
        \partsitem \Sage{} memperoleh $s\approx 25\,857$ dan $a\approx 30\,596$.
\begin{lstlisting}
sage: system = [(20053/25448)*s -  (2664/30346)*a == 17689,
....:             -(48/25448)*s + (21316/30346)*a == 21443]
sage: solve(system, s,a)
[[s == (2762271457112/106830469), a == (6537219545323/213660938)]]
sage: n(2762271457112/106830469), n(6537219545323/213660938)
(25856.5883213711, 30596.2316112410)
\end{lstlisting}
        \partsitem \Sage{} memperoleh $s\approx 25\,984$ dan $a\approx 30\,597$.
\begin{lstlisting}
sage: system = [(20053/25448)*s -  (2664/30346)*a == 17789,
....:             -(48/25448)*s + (21316/30346)*a == 21443]
sage: solve(system, s,a)
[[s == (2775832696312/106830469), a == (6537292375723/213660938)]]
sage: n(2775832696312/106830469), n(6537292375723/213660938)
(25983.5300012771, 30596.5724802865)          
\end{lstlisting}
      \end{exparts}
    \end{answer}
  \item
    Untuk sistem baja-otomotif,
    dengan perkiraan permintaan eksternal $17,589$ dan $21,243$
    yang dibahas di atas, berapakah nilai baja yang digunakan oleh industri
    baja, nilai baja yang digunakan oleh industri otomotif, dan seterusnya?
    \begin{answer}
      Sesi \textit{Sage} ini
\begin{lstlisting}
sage: var('a,s')
(a, s)
sage: eqns=[(20053/25448)*s - (2664/30346)*a == 17589,
....:       (-48/25448)*s + (21316/30346)*a == 21243] 
sage: solve(eqns, s, a)
[[s == (2745320544312/106830469), a == (6476293881123/213660938)]]
sage: s_by_s = (5395/25448)*2745320544312/106830469
sage: s_by_s
582010544505/106830469
sage: n(s_by_s)
5447.98267716114
sage: s_by_a = (2664/30346)*6476293881123/213660938  
sage: n(s_by_a)
2660.93449955743
sage: a_by_s = (48/25448)*2745320544312/106830469
sage: n(a_by_s)
48.4713936058822
sage: a_by_a = (9030/30346)*6476293881123/213660938
sage: n(a_by_a)
9019.60905818451        
\end{lstlisting}
      menghasilkan tabel berikut.
      \begin{center}
        \begin{tabular}{r|rrrr}
         \multicolumn{1}{c}{\ }  %put the | in the right place
         &\multicolumn{1}{r}{\begin{tabular}[b]{@{}c@{}}
                                \textit{digunakan oleh} \\[-.65ex] \textit{baja}
                             \end{tabular}}
         &\multicolumn{1}{r}{\begin{tabular}[b]{@{}c@{}}
                                \textit{digunakan oleh} \\[-.65ex] \textit{otomotif}
                             \end{tabular}}
         &\multicolumn{1}{r}{\begin{tabular}[b]{@{}c@{}}
                                \textit{digunakan oleh} \\[-.65ex] \textit{lainnya}
                              \end{tabular}}
         &\textit{total}                                                \\
         \cline{2-5}
         \begin{tabular}{r} \textit{nilai} \\[-.65ex] \textit{baja} \end{tabular}
           &$5\,448$  &$2\,661$  &$17\,589$     &$25\,698$                          \\
         \begin{tabular}{r} \textit{nilai} \\[-.65ex] \textit{mobil} \end{tabular}
           &$48$      &$9\,020$  &$21\,243$     &$30\,311$
        \end{tabular}
      \end{center}
      Sebagai pembanding, berikut tabel semula.
      \begin{center}
        \begin{tabular}{r|rrrr}
         \multicolumn{1}{c}{\ }  %put the | in the right place
         &\multicolumn{1}{r}{\begin{tabular}[b]{@{}c@{}}
                                \textit{digunakan oleh} \\[-.65ex] \textit{baja}
                             \end{tabular}}
         &\multicolumn{1}{r}{\begin{tabular}[b]{@{}c@{}}
                                \textit{digunakan oleh} \\[-.65ex] \textit{otomotif}
                             \end{tabular}}
         &\multicolumn{1}{r}{\begin{tabular}[b]{@{}c@{}}
                                \textit{digunakan oleh} \\[-.65ex] \textit{lainnya}
                              \end{tabular}}
         &\textit{total}                                                \\
         \cline{2-5}
         \begin{tabular}{r} \textit{nilai} \\[-.65ex] \textit{baja} \end{tabular}
           &$5\,395$  &$2\,664$  &$17\,389$     &$25\,448$                          \\
         \begin{tabular}{r} \textit{nilai} \\[-.65ex] \textit{mobil} \end{tabular}
           &$48$      &$9\,030$  &$21\,268$     &$30\,346$
        \end{tabular}
      \end{center}
    \end{answer}
  \item
    Dalam sistem baja-otomotif, rasio penggunaan baja oleh industri
    otomotif adalah $2,664/30,346$, yakni sekitar $0.0878$.
    Bayangkan suatu proses baru untuk membuat mobil mengurangi
    rasio tersebut menjadi $.0500$.
    \begin{exparts}
      \partsitem Bagaimana prediksi produksi total tahun depan akan
        berubah dibandingkan dengan contoh pertama yang dibahas
        di atas (yakni, dengan menetapkan permintaan eksternal tahun depan
        sebesar $17,589$ untuk baja dan $21,243$ untuk mobil)?
      \partsitem Prediksikan total tahun depan jika, selain itu,
        permintaan eksternal terhadap
        mobil naik menjadi $21,500$ karena mobil baru lebih murah.
    \end{exparts}
    \begin{answer}
      \begin{exparts}
        \partsitem \Sage{} memberikan $s=24,244$, $a=30,307$.
\begin{lstlisting}
 sage: var('s,a')
(s, a)
sage: system = [(20053/25448)*s -  0.0500*a == 17589,
....:             -(48/25448)*s + (21316/30346)*a == 21243]
sage: solve(system, s,a)
[[s == (12951731858499/534221147), a == (32381469405615/1068442294)]]
sage: n(12951731858499/534221147), n(32381469405615/1068442294)
(24244.1392131918, 30307.1767071166)         
\end{lstlisting}
        \partsitem Sage memberikan $s=24,267$, $a=30,673$.
\begin{lstlisting}
sage: system = [(20053/25448)*s -  0.0500*a == 17589,
....:             -(48/25448)*s + (21316/30346)*a == 21500]
sage: solve(system, s,a)
[[s == (12964136043940/534221147), a == (16386224431390/534221147)]]
sage: n(12964136043940/534221147), n(16386224431390/534221147)
(24267.3584090448, 30673.1107957993)          
\end{lstlisting}
      \end{exparts}
    \end{answer}
  \item
    Tabel ini memberikan angka untuk sistem otomotif-baja dari
    tahun yang berbeda, yakni 1947 (lihat \cite{Leontief1951}).
    Satuannya adalah miliaran dolar tahun 1947.
    \begin{center}
      \begin{tabular}{r|rrrr}
             \multicolumn{1}{c}{\ }  % omit the | 
             &\multicolumn{1}{r}{\begin{tabular}[b]{@{}c@{}}
                                   \textit{digunakan oleh} \\[-.65ex] \textit{baja}
                                 \end{tabular}}
             &\multicolumn{1}{r}{\begin{tabular}[b]{@{}c@{}}
                                   \textit{digunakan oleh} \\[-.65ex] \textit{otomotif}
                                 \end{tabular}}
             &\multicolumn{1}{r}{\begin{tabular}[b]{@{}c@{}}
                                   \textit{digunakan oleh} \\[-.65ex] \textit{lainnya}
                                 \end{tabular}}
             &\multicolumn{1}{r}{\textit{total}}             \\ \cline{2-5}
        \begin{tabular}{r} \textit{nilai} \\[-.65ex]
                           \textit{baja}
            \end{tabular}
            &$6.90$  &$1.28$  &   &$18.69$                              \\
        \begin{tabular}{r} \textit{nilai} \\[-.65ex] \textit{mobil}
            \end{tabular}
            &$0$     &$4.40$  &   &$14.27$
      \end{tabular}
    \end{center}
    \begin{exparts}
      \partsitem Hitung keluaran total jika permintaan eksternal
         tahun depan adalah: permintaan baja naik $10$\% dan permintaan
         mobil naik $15$\%.
      \partsitem Bagaimana perbandingan rasio-rasio ini dengan rasio yang
         diberikan di atas dalam pembahasan perekonomian tahun 1958?
      \partsitem Selesaikan persamaan tahun 1947 dengan permintaan eksternal
         tahun 1958 (perhatikan perbedaan satuannya; satu dolar tahun 1947
         dapat membeli kira-kira sebanyak yang dapat dibeli oleh
         \$$1.30$ pada tahun 1958).
         Seberapa jauh prediksi keluaran totalnya meleset?
    \end{exparts}
    \begin{answer}
      \begin{exparts}
        \partsitem
          Berikut persamaannya.
	  \begin{equation*}
	    \begin{linsys}{2}
	       (11.79/18.69)s   &-   &(1.28/14.27)a &= &10.51 \\
	       -(0/18.69)s   &+   &(9.87/14.27)a &= &9.87
	    \end{linsys}
	  \end{equation*}
          \textit{Sage} memberikan hasil ini (termasuk pemeriksaan awal).
\begin{lstlisting}
sage: var('a,s')
(a, s)
sage: eqns=[(11.79/18.69)*s - (1.28/14.27)*a == 10.51,
....:       (-0)*s         + (9.87/14.27)*a == 9.87]
sage: solve(eqns,s,a)
[[s == (1869/100), a == (1427/100)]]
sage: n(1869/100)
18.6900000000000
sage: n(1427/100)
14.2700000000000
sage: eqns=[(11.79/18.69)*s - (1.28/14.27)*a == 11.561,
....:       (-0)*s         + (9.87/14.27)*a == 11.3505]
sage: solve(eqns,s,a)
[[s == (8119559/393000), a == (32821/2000)]]
sage: n(8119559/393000)
20.6604554707379
sage: n(32821/2000)
16.4105000000000            
\end{lstlisting}
	\partsitem
          Berikut rasio tahun 1947, hingga tiga tempat desimal.
	  \begin{center}
            \begin{tabular}{r|cc}
               1947   &\textit{oleh baja}  &\textit{oleh otomotif}  \\ \hline
               \textit{penggunaan baja}  &$(6.90/18.69)=0.369$ &$(1.28/14.27)=0.090$  \\
               \textit{penggunaan mobil}  &$(0/18.69)=0.00$ &$(4.40/14.27)=0.308$
            \end{tabular}
          \end{center}
          Berikut rasio tahun 1958.
          \begin{center}
            \begin{tabular}{r|cc}
               1958           &\textit{oleh baja}  &\textit{oleh otomotif}  \\ \hline
               \textit{penggunaan baja}  &$(5395/25448)=0.212$ &$(2664/30346)=0.088$  \\
               \textit{penggunaan mobil}  &$(48/25448)=0.002$ &$(9030/30346)=0.298$
            \end{tabular}
          \end{center}
	\partsitem
          Permintaan eksternal pada tahun 1958 adalah $17.589$ dan~$21.243$,
          dalam miliaran dolar tahun 1958.
          Dalam dolar tahun 1947, nilainya adalah $(17.589/1.3)=13.53$ dan
          $(21.243/1.3)=16.34$.
          \textit{Sage} memberikan angka-angka berikut.
\begin{lstlisting}
sage: eqns=[(11.79/18.69)*s - (1.28/14.27)*a == 13.53,
....:       (-0)*s         + (9.87/14.27)*a == 16.34]
sage: solve(eqns,s,a)
[[s == (137466107/5541300), a == (1165859/49350)]]
sage: n(137466107/5541300)
24.8075554472777
sage: n(1165859/49350)
23.6242958459980            
\end{lstlisting}
      Dengan mengonversi kembali ke dolar tahun 1958 melalui perkalian dengan
      $1.3$, kita memperoleh $32.25$ dan~$30.71$.
      Dibandingkan dengan prediksi berdasarkan data tahun 1958, yakni sekitar
      $25.70$ dan~$30.31$, hasil tersebut lebih tinggi sekitar $6.55$ dan
      $0.40$, atau masing-masing sekitar $25.5$\% dan $1.3$\%.
      \end{exparts}
    \end{answer}
  \item
    Prediksikan produksi total setiap sektor pada tahun depan untuk ketiga
    sektor perekonomian hipotetis yang ditunjukkan di bawah ini
    \begin{center}
      \begin{tabular}{r|rrrrr}
             \multicolumn{1}{c}{\ }  % omit the |
             &\multicolumn{1}{r}{\begin{tabular}[b]{@{}c@{}}
                                     \textit{digunakan oleh} \\[-.65ex] \textit{pertanian}
                                 \end{tabular}}
             &\multicolumn{1}{r}{\begin{tabular}[b]{@{}c@{}}
                                    \textit{digunakan oleh} \\[-.65ex] \textit{kereta api}
                                 \end{tabular}}
             &\multicolumn{1}{r}{\begin{tabular}[b]{@{}c@{}}
                                    \textit{digunakan oleh} \\[-.65ex]
                                    \textit{pelayaran}
                                 \end{tabular}}
             &\multicolumn{1}{r}{\begin{tabular}[b]{@{}c@{}}
                                   \textit{digunakan oleh} \\[-.65ex] \textit{lainnya}
                                 \end{tabular}}
             &\textit{total}                                                \\
             \cline{2-6}
        \begin{tabular}{r} \textit{nilai} \\[-.65ex] \textit{pertanian}
             \end{tabular}
             &$25$  &$50$  &$100$ &     &$800$             \\
        \begin{tabular}{r} \textit{nilai} \\[-.65ex] \textit{kereta api}
             \end{tabular}
             &$25$  &$50$  &$50$  &     &$300$             \\
        \begin{tabular}{r} \textit{nilai} \\[-.65ex] \textit{pelayaran}
             \end{tabular}
             &$15$  &$10$  &$0$   &     &$500$
      \end{tabular}
    \end{center}
    jika permintaan eksternal tahun depan adalah sebagaimana dinyatakan.
    \begin{exparts}
      \partsitem $625$ untuk pertanian, $200$ untuk kereta api, $475$ untuk pelayaran
      \partsitem $650$ untuk pertanian, $150$ untuk kereta api, $450$ untuk pelayaran
    \end{exparts}
    % ID-EDITION-SUPPLIED-ANSWER: not present in the pinned authority source
    % or in the official jhanswer.pdf; independently derived and verified.
    \begin{answer}
      \textit{Catatan edisi Indonesia: buku jawaban resmi tidak memuat
      jawaban untuk latihan ini.  Penyelesaian berikut diturunkan secara
      independen dari model di atas.}
      Misalkan $f$, $r$, dan $p$ berturut-turut menyatakan produksi total
      pertanian, kereta api, dan pelayaran tahun depan, dan misalkan
      $d_f$, $d_r$, dan $d_p$ menyatakan permintaan eksternalnya.  Rasio
      masukan terhadap keluaran memberikan sistem
      \begin{equation*}
        \begin{linsys}{3}
          (31/32)f  &- &(1/6)r  &- &(1/5)p  &= &d_f \\
          -(1/32)f  &+ &(5/6)r  &- &(1/10)p &= &d_r \\
          -(3/160)f &- &(1/30)r &+ &p       &= &d_p.
        \end{linsys}
      \end{equation*}
      \begin{exparts}
        \partsitem Untuk $(d_f,d_r,d_p)=(625,200,475)$, penyelesaiannya
          adalah
          \begin{align*}
            f&=\frac{3\,074\,400}{3\,817}\approx805{,}449, &
            r&=\frac{1\,260\,900}{3\,817}\approx330{,}338, &
            p&=\frac{1\,912\,750}{3\,817}\approx501{,}113.
          \end{align*}
        \partsitem Untuk $(d_f,d_r,d_p)=(650,150,450)$, penyelesaiannya
          adalah
          \begin{align*}
            f&=\frac{3\,110\,400}{3\,817}\approx814{,}881, &
            r&=\frac{1\,020\,900}{3\,817}\approx267{,}461, &
            p&=\frac{1\,810\,000}{3\,817}\approx474{,}194.
          \end{align*}
      \end{exparts}
    \end{answer}
  \item
      Tabel ini menunjukkan hubungan di antara tiga segmen suatu
      perekonomian (lihat \cite{BangorRpt}).
      \begin{center}
        \begin{tabular}{r|rrrrr}
             \multicolumn{1}{c}{\ }  % omit the |
             &\multicolumn{1}{r}{\begin{tabular}[b]{@{}c@{}}
                                    \textit{digunakan oleh} \\[-.65ex] \textit{pangan}
                                 \end{tabular}}
             &\multicolumn{1}{r}{\begin{tabular}[b]{@{}c@{}}
                                   \textit{digunakan oleh} \\[-.65ex]
                                   \textit{grosir}
                                 \end{tabular}}
             &\multicolumn{1}{r}{\begin{tabular}[b]{@{}c@{}}
                                    \textit{digunakan oleh} \\[-.65ex] \textit{eceran}
                                 \end{tabular}}
             &\multicolumn{1}{r}{\begin{tabular}[b]{@{}c@{}}
                                    \textit{digunakan oleh} \\[-.65ex] \textit{lainnya}
                                 \end{tabular}}
             &\textit{total}                                             \\
          \cline{2-6}
          \begin{tabular}[b]{r} \textit{nilai} \\[-.65ex] \textit{pangan}
              \end{tabular}
              &$0$    &$2\,318$   &$4\,679$   &     &$11\,869$   \\
          \begin{tabular}[b]{r} \textit{nilai} \\[-.65ex] \textit{grosir}
              \end{tabular}
              &$393$  &$1\,089$   &$22\,459$  &     &$122\,242$   \\
          \begin{tabular}[b]{r} \textit{nilai} \\[-.65ex] \textit{eceran}
              \end{tabular}
              &$3$    &$53$       &$75$       &     &$116\,041$   \\
        \end{tabular}
      \end{center}
      Kita akan melakukan Analisis Masukan-Keluaran pada
      sistem ini.
      \begin{exparts}
        \partsitem Lengkapi angka permintaan eksternal
          tahun ini.
        \partsitem Susun sistem linearnya dengan membiarkan permintaan
          eksternal tahun depan kosong.
        \partsitem Selesaikan sistem ketika permintaan eksternal tahun depan
          diperoleh dengan menaikkan permintaan eksternal tahun ini
          sebesar $10$\%.
          Apakah kegiatan usaha total ketiga sektor meningkat sebesar
          $10$\%?
          Apakah ketiganya bahkan meningkat dengan laju yang sama?
        \partsitem Selesaikan sistem ketika permintaan eksternal tahun depan
          diperoleh dengan menurunkan permintaan eksternal tahun ini
          sebesar $7$\%.
          (Studi yang menjadi sumber angka-angka ini menyimpulkan bahwa
          akibat penutupan fasilitas militer setempat, pendapatan pribadi
          keseluruhan di wilayah tersebut akan turun $7$\%, sehingga ini
          dapat menjadi perkiraan awal tentang apa yang sebenarnya akan terjadi.)
      \end{exparts}
      % ID-EDITION-SUPPLIED-ANSWER: not present in the pinned authority source
      % or in the official jhanswer.pdf; independently derived and verified.
      \begin{answer}
        \textit{Catatan edisi Indonesia: buku jawaban resmi tidak memuat
        jawaban untuk latihan ini.  Penyelesaian berikut diturunkan secara
        independen dari model di atas.}
        Misalkan $p$, $g$, dan $e$ berturut-turut menyatakan produksi total
        sektor pangan, grosir, dan eceran tahun depan.
        \begin{exparts}
          \partsitem Permintaan eksternal tahun ini adalah
            \begin{align*}
              d_p&=11\,869-0-2\,318-4\,679=4\,872,\\
              d_g&=122\,242-393-1\,089-22\,459=98\,301,\\
              d_e&=116\,041-3-53-75=115\,910.
            \end{align*}
          \partsitem Dengan $d_p$, $d_g$, dan $d_e$ sebagai permintaan
            eksternal tahun depan, sistem linearnya adalah
            \begin{equation*}
              \begin{linsys}{3}
                p              &- &(2\,318/122\,242)g
                                  &- &(4\,679/116\,041)e &= &d_p \\
                -(393/11\,869)p &+ &(121\,153/122\,242)g
                                  &- &(22\,459/116\,041)e &= &d_g \\
                -(3/11\,869)p   &- &(53/122\,242)g
                                  &+ &(115\,966/116\,041)e &= &d_e.
              \end{linsys}
            \end{equation*}
          \partsitem Kenaikan $10\%$ memberikan permintaan eksternal
            $(5\,359{,}2,\ 108\,131{,}1,\ 127\,501)$ dan penyelesaian
            \[
              (p,g,e)=(13\,055{,}9,\ 134\,466{,}2,\ 127\,645{,}1).
            \]
            Ya.  Produksi total ketiga sektor meningkat tepat $10\%$;
            ketiganya meningkat dengan laju yang sama.
          \partsitem Penurunan $7\%$ memberikan permintaan eksternal
            $(4\,530{,}96,\ 91\,419{,}93,\ 107\,796{,}30)$ dan penyelesaian
            \[
              (p,g,e)=(11\,038{,}17,\ 113\,685{,}06,\ 107\,918{,}13).
            \]
            Dalam hal ini produksi total ketiga sektor juga turun tepat
            $7\%$.
        \end{exparts}
      \end{answer}
\end{exercises}
\index{Analisis Masukan-Keluaran|)}
\endinput
